% Capitulo: A estrutura dos times de dados
\chapter{A estrutura dos times de dados}
\label{cap:estrutura-times}

Antes de falar de como um time de dados trabalha, vale olhar para quem compõe esse time. A
estrutura mais comum nas empresas médias e grandes se formou ao longo das ondas descritas na
Parte~\ref{parte:origens}, e cada onda acrescentou um papel. O resultado é um time grande,
especializado e com muitas passagens de bastão. Este capítulo descreve essa estrutura, mostra
o que a IA comprime nela e o que não comprime, propõe uma versão enxuta possível e termina
com algumas apostas sobre o futuro.

\section{A estrutura tradicional}
\label{sec:estrutura-tradicional}

Numa estrutura típica, há um responsável pela área de dados, abaixo dele um gerente de
engenharia, com um time de engenheiros de dados, e um gerente de analytics, com um time de
analistas de dados e de BI. Muitas vezes há também um product manager, que traduz as
demandas do negócio em prioridades, e alguém dedicado a práticas ágeis, como um scrum master
ou um agile coach. E isso sem contar os times de negócio, que têm os seus próprios analistas e
gestores e são, ao mesmo tempo, clientes e fontes de conhecimento para o time de dados
(parte de cima da Figura~\ref{fig:estrutura-times}).

Essa estrutura surgiu por bons motivos. As ferramentas de cada etapa eram especializadas e
difíceis: extrair e carregar dados exigia uma formação, modelar exigia outra, montar painéis
exigia outra. Ter gente dedicada a cada etapa era a forma de ter qualidade em cada uma.

O custo aparece no caminho de um pedido. Uma pergunta de negócio chega ao product manager,
vira uma história no quadro, vai para a engenharia, que prepara os dados, volta para os
analistas, que montam o painel, e só então chega a quem perguntou. Cada passagem perde um
pouco de contexto e acrescenta um pouco de espera. E, quando o número sai errado, cada etapa
tende a achar que o problema está na outra.

Há também um custo menos visível, conhecido há meio século. Frederick Brooks observou que o
número de canais de comunicação num time cresce muito mais rápido do que o número de pessoas:
com $n$ pessoas, há $n(n-1)/2$ pares que podem precisar conversar \parencite{brooks1975}. Um
time de quatro pessoas tem seis canais; um de doze tem sessenta e seis. Parte do trabalho de
um time grande é coordenar a si mesmo.

E há a observação de Melvin Conway, de 1968: a estrutura dos sistemas que uma organização
constrói tende a copiar a estrutura de comunicação dessa organização \parencite{conway1968}.
Em dados, dois times separados, um de engenharia e um de visualização, costumam produzir duas
camadas que não conversam bem, exatamente como o capítulo~\ref{cap:pipeline-dashboard}
descreveu.

\begin{figure}[htbp]
  \centering
  % Estrutura tradicional de um time de dados e uma versao enxuta possivel (cap. 23).
\begin{tikzpicture}[diagrama, x=1mm, y=1mm,
    titulo/.style={anchor=west, font=\footnotesize\scshape, text=cinza},
    cargo/.style={caixa, minimum width=24mm, minimum height=8mm, font=\scriptsize},
    pilha/.style={cargo, double copy shadow={shadow xshift=1mm, shadow yshift=-1mm, draw=tinta, fill=white}},
    negocio/.style={caixa, draw=cinza, dashed, font=\scriptsize, text=cinza},
    hierarquia/.style={draw=tinta, line width=0.5pt}]

  % --- estrutura tradicional
  \node[titulo] at (0,58) {estrutura tradicional};
  \node[cargo] (head) at (58,50) {Head de dados};
  \node[cargo] (ge) at (14,37) {Gerente de\\engenharia};
  \node[cargo] (ga) at (44,37) {Gerente de\\analytics};
  \node[cargo] (pm) at (74,37) {Product\\manager};
  \node[cargo] (ag) at (103,37) {Agilista\\(scrum master)};
  \foreach \n in {ge,ga,pm,ag} \draw[hierarquia] (head.south) -- ++(0,-3) -| (\n.north);
  \node[pilha] (eng) at (14,23) {Engenheiros\\de dados};
  \node[pilha] (ana) at (44,23) {Analistas de\\dados e BI};
  \draw[hierarquia] (ge) -- (eng);
  \draw[hierarquia] (ga) -- (ana);
  \node[negocio, minimum width=50mm, minimum height=8mm] (neg1) at (88,23)
    {Times de negócio: analistas de\\negócio, gestores, usuários};
  \draw[ligacao, dashed] (pm.south) -- (pm.south |- neg1.north);

  % --- versao enxuta
  \draw[fio] (0,12) -- (118,12);
  \node[titulo] at (0,7) {uma versão enxuta possível};
  \node[destaquecaixa, minimum width=34mm, minimum height=8mm, font=\scriptsize] (lider) at (58,0)
    {Líder de dados\\(produto e fluxo)};
  \node[pilha] (prof) at (40,-14) {Profissionais de\\dados (ponta a ponta)};
  \node[caixa, draw=cinza, minimum width=24mm, minimum height=8mm, font=\scriptsize] (ia) at (76,-14)
    {Agentes de IA\\(revisados pelo time)};
  \draw[hierarquia] (lider.south) -- ++(0,-3) -| (prof.north);
  \draw[ligacao] (lider.south) -- ++(0,-3) -| (ia.north);
  \node[negocio, minimum width=26mm, minimum height=8mm] (neg2) at (105,-14)
    {Negócio, com\\analistas próprios};
  \draw[ligacao, dashed] (ia.east) -- (neg2.west);
\end{tikzpicture}

  \caption{A estrutura tradicional de um time de dados e uma versão enxuta possível com
    a IA. Em linha tracejada, a relação com o negócio.}
  \label{fig:estrutura-times}
  \fonte{elaborado pelo autor.}
\end{figure}

\rascunho{
  Relato: as estruturas de time que eu vi ou montei, quantas passagens de bastão havia entre
  o pedido e a entrega, e o que isso custava. Sem nomes.
}

\section{O que a IA comprime}
\label{sec:ia-comprime}

A IA não elimina etapas do trabalho com dados, mas encurta muitas delas e, principalmente,
diminui a distância entre elas. Com um assistente ajudando nas partes que não domina, a mesma
pessoa consegue ir da ingestão ao painel: um analista que sabe SQL e negócio escreve o modelo
com testes, um engenheiro monta um painel razoável, um analytics engineer cuida do caminho
inteiro (\veja{sec:elt}). A fronteira entre engenharia e análise, que justificava dois times,
fica mais fina.

Alguns papéis também tendem a se juntar. Em times menores, o líder de dados acumula a gestão
de produto e a gestão do fluxo, e as práticas ágeis passam a ser responsabilidade do time,
não de um cargo dedicado. Não porque esses papéis tenham perdido importância, mas porque um
time de quatro pessoas raramente justifica uma pessoa só para cada um deles.

O que a IA não comprime é igualmente importante. Conhecer o negócio continua exigindo tempo e
contato com as pessoas (capítulo~\ref{cap:negocio}). Revisar o que a IA produz continua
exigindo quem saiba ler código e conferir números (capítulo~\ref{cap:programar}). E decidir o
que construir, com que frequência e a que custo continua sendo trabalho de gente
(capítulos~\ref{cap:frequencia} e~\ref{cap:custo-consumo}). Um time enxuto que perde essas
capacidades não ficou eficiente; ficou frágil.

Há dois riscos a evitar. O primeiro é reduzir a estrutura sem reduzir a demanda: um time menor
com o mesmo volume de pedidos e a mesma forma de trabalhar vira um time sobrecarregado, e a
qualidade cai justamente na revisão, que é o que a IA mais exige. O segundo é cortar as
posições de entrada. Se o trabalho de júnior encolher e ninguém mais for contratado para
aprender, daqui a alguns anos faltarão os seniores que hoje revisam o que a IA faz.

\section{Uma estrutura enxuta}
\label{sec:estrutura-enxuta}

A parte de baixo da Figura~\ref{fig:estrutura-times} mostra uma possibilidade. Um líder de
dados, que cuida do produto e do fluxo de trabalho. Alguns profissionais de dados que
trabalham de ponta a ponta, da fonte ao painel, cada um com uma especialidade mais funda, mas
todos capazes de atravessar as camadas. Agentes de IA executando parte das tarefas, sempre
revisados pelo time. E o negócio mais perto dos dados, com analistas próprios que consomem as
bases certificadas e fazem as suas análises (\veja{sec:governanca}).

Duas referências ajudam a pensar essa organização com mais cuidado. \textcite{skelton2019}
propõem um vocabulário para tipos de time: times alinhados a um fluxo de valor, que entregam
ao negócio; times de plataforma, que oferecem serviços internos para os outros; e times de
apoio, que ajudam os demais a adquirir uma capacidade. E propõem formas de interação entre
eles, da colaboração próxima ao consumo de um serviço. \textcite{dehghani2022} vai além e
propõe que os próprios domínios de negócio sejam donos dos seus produtos de dados, com uma
plataforma comum de autoatendimento e uma governança federada.

Nenhuma das duas propostas precisa ser adotada por inteiro. O que vale extrair delas é a
pergunta: o que fica centralizado numa plataforma, com poucas pessoas e muita automação, e o
que vai para perto do negócio? A IA torna viável centralizar menos, porque cada pessoa
consegue fazer mais sozinha, mas não dispensa a plataforma, a governança e as definições
comuns que impedem cada área de criar a sua própria verdade.

O Scrum com Kanban descrito nos próximos capítulos foi pensado para funcionar também num time
assim, de três ou quatro pessoas, sem cargos dedicados a cada prática.

\section{Uma previsão para o futuro}
\label{sec:previsao-times}

Previsões sobre tecnologia envelhecem mal, e esta deve ser lida como aposta. Algumas
tendências, porém, já são visíveis o bastante para arriscar uma direção.

Os times de dados tendem a ficar menores e mais seniores, com a IA fazendo parte do trabalho
que antes era de pessoas em início de carreira: escrever a primeira versão de uma consulta,
documentar tabelas, montar painéis padronizados. O valor de cada profissional passa a estar
mais no julgamento do que no volume do que produz.

O analista de dados e o analista de negócio tendem a convergir. Se a técnica fica mais
acessível, a pessoa que conhece o negócio consegue fazer mais análise sozinha, e a que sabe
análise precisa conhecer mais o negócio para se diferenciar. O ponto de encontro é um
profissional que vive perto da área de negócio e domina o suficiente da técnica para
trabalhar com a IA com critério.

O gerente de dados tende a acumular três gestões: a do produto (o que construir), a do fluxo
(como o trabalho anda) e a do consumo (quanto custa, em créditos e tokens). E, com a IA
proativa do capítulo~\ref{cap:estagios-ia}, uma quarta: a das regras e alçadas que dizem o
que a IA pode decidir sozinha.

E há uma pergunta em aberto, que nenhuma empresa respondeu bem ainda: como se forma um sênior
se o trabalho de júnior encolher? A resposta provavelmente passa por usar a IA como tutor, e
não como substituto, na formação de quem entra (\veja{sec:ia-tutor}), mas ainda está para ser
construída.

\rascunho{
  A minha aposta pessoal, em primeira pessoa: qual destas tendências eu acho mais provável,
  qual eu acho exagerada, e o que eu estou fazendo no meu próprio time por causa disso.
  Alinhar com o epílogo para não repetir.
}

\begin{leituras}
  \leitura{skelton2019}{Um vocabulário claro para tipos de time (de fluxo, de plataforma,
    de apoio) e para as formas de interação entre eles.}
  \leitura{dehghani2022}{A proposta de organizar dados por domínio de negócio, com times
    responsáveis pelos próprios produtos de dados.}
  \leitura{brooks1975}{Um clássico sobre por que times maiores nem sempre entregam mais. A
    observação sobre os canais de comunicação continua atual.}
\end{leituras}
